\documentclass[../../main.tex]{subfiles}
\begin{document}
\begin{oframed}
  \subfile{problem}
\end{oframed}

{\bf [解答]}
\paragraph{(1)}
$f'(x)=2ax+b$である．題意より$x=0,1,-1$で$|f(x)|\le 1$だから
\begin{align}
  \begin{dcases}
    |f(0)| = |c| \le 1 \\
    |f(1)| = |a+b+c| \le 1 \\
    |f(-1)| = |a-b+c| \le 1
  \end{dcases}
\end{align}
となる．絶対値を外して
\begin{align}
  -1 \le c \le 1 \label{eq:1}\\
  -1 \le a+b+c \le 1 \label{eq:2}\\
  -1 \le a-b+c \le 1 \label{eq:3}
\end{align}
である．ここで
\begin{align}
  \dfrac{3}{2}\text{\cref{eq:2}} + \dfrac{1}{2}\text{\cref{eq:3}} - 2\text{\cref{eq:1}} \label{eq:4}
\end{align}
を計算すると
\begin{align}
  -4 \le 2a + b \le 4
\end{align}
である．$f'(1)=2a+b$だから
\begin{align}
  |f'(1)| \le 4
\end{align}
であり，よって題意は示された．

% 元の解
% $g(x)=ax^2+bx$とおくと，$f'(x)=g'(x)=2ax+b$である．又，題意から
% \begin{align}
% -2\le g(1)-g(0)\le2, \qquad -2\le g(-1)-g(0)\le2 \iff -2\le\pm a+b\le2 \quad\cdots\text{①}
% \end{align}
% が成立する．さらに，$g'(1)=2a+b$である．ここで，題意が成立しないと仮定すると，
% \begin{align}
% |2a+b|>4 \quad\cdots\text{②}
% \end{align}
% ところが，①$\wedge$②をみたす$(a,b)$は存在せず，不適．よって背理法から題意は成立．

\paragraph{(2)}

$|f'(1)|=4$が成り立つということは(1)より\cref{eq:4}で等号が成立するので\cref{eq:1,eq:2,eq:3}より
\begin{align}
  c = \mp 1 \label{eq:5}\\
  a+b+c = \pm 1 \label{eq:6}\\
  a-b+c = \pm 1 \label{eq:7}
\end{align}
である．ただし複合同順である．\cref{eq:6,eq:7}を辺々足し引きして
\begin{align}
  \begin{dcases}
    2a+2c = \pm 2 \quad \therefore a+c = \pm 1 \\
    2b = 0 \quad \therefore b = 0
  \end{dcases}\label{eq:8}
\end{align}
である．さらに\cref{eq:5}を\cref{eq:8}の第一式に代入して
\begin{align}
  a = \pm 2
\end{align}
である．以上より
\begin{align}
  (a,b,c) = (\pm 2,0,\mp 1)
\end{align}
である．ただし複合同順である．この時関数$f(x)$は
\begin{align}
  f(x) = \pm \left(2x^2-1\right)
\end{align}
であるが，これは題意の条件$|f(x)|\le 1\, (|x|\le 1)$を満たすから十分である．
以上より求める関数は
\begin{align}
  f(x) = \pm \left(2x^2-1\right)
\end{align}
である．

{\bf [別解]}

[解1]では不等式の性質を使って$x=0,\pm 1$の値のみを使って解いたが，もっとゴリゴリに二次方程式の性質を使って解くこともできるのでその例も示そう．受験だとこういう解答の方が途中で事故りにくいので安全かもしれない．

\paragraph{(1)}
$a=0$での成立は明らかである（$f(x)=bx+c$より$|b|\le 1$となり，$|f'(1)|=|b|\le 1<4$）．また$f(x)$を$-f(x)$に置き換えても条件$|f(x)|\le 1\, (|x|\le 1)$と評価対象$|f'(1)|$は変わらないため，対称性から以下$a>0$としてよい．$f(x)$は二次関数だから平方完成すると
\begin{align}
  f(x)=a\left(x+\dfrac{b}{2a}\right)^2-\dfrac{b^2}{4a}+c
\end{align}
であり，区間内での最大最小値は軸の位置$x=-\dfrac{b}{2a}$の値によって以下のようになる．
\begin{table}[htb]
  \centering
  \caption{$f(x)$の軸の位置と，対応する$f(x)$の最大最小値}
  \begin{tabular}{c|c|c|c}
    & $1^{\circ} \ $ $-1\le-\dfrac{b}{2a}\le0$ & $2^{\circ} \ $ $0\le-\dfrac{b}{2a}\le1$ & $3^{\circ} \ $ otherwise \\
    \hline
    $\max f$ & $f(1)=a+b+c$ & $f(-1)=a-b+c$ & $f(\pm 1)$のうち小さくない方 \\
    \hline
    $\min f$ & $f\left(-\dfrac{b}{2a}\right)=-\dfrac{b^2}{4a}+c$ & $f\left(-\dfrac{b}{2a}\right)=-\dfrac{b^2}{4a}+c$ & $f(\pm 1)$のうち大きくない方
  \end{tabular}
\end{table}

\begin{figure}[htb]
  \centering
  \begin{tikzpicture}[x=1.4cm, y=1.3cm, >=stealth, every node/.style={font=\small}]
    % --- 左：$1^\circ$ (-1 <= 軸 <= 0) ---
    \begin{scope}[xshift=0cm]
      \node[font=\footnotesize\bfseries] at (0, 1.65) {$1^\circ\ (-1 \le -\frac{b}{2a} \le 0)$};
      \draw[->, thick, gray!80!black] (-1.4, 0) -- (1.4, 0) node[right, font=\scriptsize] {$x$};
      \draw[->, thick, gray!80!black] (0, -1.15) -- (0, 1.15) node[above, font=\scriptsize] {$y$};
      \draw[densely dotted, thick, gray!70!black] (-1, -1.0) -- (-1, 1.0);
      \draw[densely dotted, thick, gray!70!black] (1, -1.0) -- (1, 1.0);
      \draw[densely dotted, thick, gray!70!black] (-0.35, 0) -- (-0.35, -0.85);
      \draw[line width=1.4pt, blue!85!black, domain=-1.2:1.15, samples=50, smooth]
      plot (\x, {0.9*(\x+0.35)*(\x+0.35)-0.85});
      \node[circle, fill=red!85!black, inner sep=1.4pt] at (-0.35, -0.85) {};
      \node[circle, fill=red!85!black, inner sep=1.4pt] at (1, 0.790) {};
      \node[below=2pt, font=\scriptsize, fill=white, inner sep=0.8pt] at (-1, 0) {$-1$};
      \node[below=2pt, font=\scriptsize, fill=white, inner sep=0.8pt] at (1, 0) {$1$};
      \node[above=2pt, font=\scriptsize] at (-0.35, 0) {$-\frac{b}{2a}$};
      \node[below=2pt, font=\scriptsize, red!85!black] at (-0.35, -0.85) {$\min$};
      \node[above left=1pt, font=\scriptsize, red!85!black] at (1, 0.790) {$\max$};
    \end{scope}

    % --- 中：$2^\circ$ (0 <= 軸 <= 1) ---
    \begin{scope}[xshift=5.6cm]
      \node[font=\footnotesize\bfseries] at (0, 1.65) {$2^\circ\ (0 \le -\frac{b}{2a} \le 1)$};
      \draw[->, thick, gray!80!black] (-1.4, 0) -- (1.4, 0) node[right, font=\scriptsize] {$x$};
      \draw[->, thick, gray!80!black] (0, -1.15) -- (0, 1.15) node[above, font=\scriptsize] {$y$};
      \draw[densely dotted, thick, gray!70!black] (-1, -1.0) -- (-1, 1.0);
      \draw[densely dotted, thick, gray!70!black] (1, -1.0) -- (1, 1.0);
      \draw[densely dotted, thick, gray!70!black] (0.35, 0) -- (0.35, -0.85);
      \draw[line width=1.4pt, blue!85!black, domain=-1.15:1.2, samples=50, smooth]
      plot (\x, {0.9*(\x-0.35)*(\x-0.35)-0.85});
      \node[circle, fill=red!85!black, inner sep=1.4pt] at (0.35, -0.85) {};
      \node[circle, fill=red!85!black, inner sep=1.4pt] at (-1, 0.790) {};
      \node[below=2pt, font=\scriptsize, fill=white, inner sep=0.8pt] at (-1, 0) {$-1$};
      \node[below=2pt, font=\scriptsize, fill=white, inner sep=0.8pt] at (1, 0) {$1$};
      \node[above=2pt, font=\scriptsize] at (0.35, 0) {$-\frac{b}{2a}$};
      \node[below=2pt, font=\scriptsize, red!85!black] at (0.35, -0.85) {$\min$};
      \node[above right=1pt, font=\scriptsize, red!85!black] at (-1, 0.790) {$\max$};
    \end{scope}

    % --- 右：$3^\circ$ (軸が区間外) ---
    \begin{scope}[xshift=11.2cm]
      \node[font=\footnotesize\bfseries] at (0, 1.65) {$3^\circ\ (\text{その他})$};
      \draw[->, thick, gray!80!black] (-1.4, 0) -- (1.45, 0) node[right, font=\scriptsize] {$x$};
      \draw[->, thick, gray!80!black] (0, -1.15) -- (0, 1.15) node[above, font=\scriptsize] {$y$};
      \draw[densely dotted, thick, gray!70!black] (-1, -1.0) -- (-1, 1.0);
      \draw[densely dotted, thick, gray!70!black] (1, -1.0) -- (1, 1.0);
      \draw[densely dotted, thick, gray!70!black] (1.28, 0) -- (1.28, -0.80);
      \draw[line width=1.4pt, blue!85!black, domain=-1.2:1.5, samples=50, smooth]
      plot (\x, {0.25*(\x-1.28)*(\x-1.28)-0.80});
      \node[circle, fill=red!85!black, inner sep=1.4pt] at (-1, 0.4996) {};
      \node[circle, fill=red!85!black, inner sep=1.4pt] at (1, -0.7804) {};
      \node[below=2pt, font=\scriptsize, fill=white, inner sep=0.8pt] at (-1, 0) {$-1$};
      \node[below=2pt, font=\scriptsize, fill=white, inner sep=0.8pt] at (1, 0) {$1$};
      \node[above=2pt, font=\scriptsize] at (1.28, 0) {$-\frac{b}{2a}$};
      \node[above right=1pt, font=\scriptsize, red!85!black] at (-1, 0.4996) {$\max$};
      \node[below left=1pt, font=\scriptsize, red!85!black] at (1, -0.7804) {$\min$};
    \end{scope}
  \end{tikzpicture}
  \caption{$a>0$ のときの軸の位置による 3 通りの場合分け（$3^\circ$ は $-\dfrac{b}{2a}>1$ の例）}
\end{figure}

\subparagraph{$1^{\circ}$の時（$-1\le-\dfrac{b}{2a}\le0$）}

$a>0$から$0\le b\le 2a$である．$|f(x)| \le 1$より$-1 \le \min f, \max f \le 1$であるから
\begin{align}
  \begin{dcases}
    -1 \le -\dfrac{b^2}{4a}+c \le 1 \\
    -1 \le a+b+c \le 1
  \end{dcases}
\end{align}
である．辺々引き算して
\begin{align}
  & -2 \le a+b+\dfrac{b^2}{4a} \le 2 \\
  & -8a \le 4a^2+4ab+b^2 \le 8a \\
  & -8a \le (2a+b)^2 \le 8a \\
  &\left|2a+b\right| \le 2\sqrt{2a} \\
  \therefore
  & -2a - 2\sqrt{2a} \le b \le -2a + 2\sqrt{2a}
\end{align}
である．これと$0\le b\le 2a$を合わせて$(a,b)$平面に図示すると以下の斜線部である．

\begin{figure}[htb]
  \centering
  \begin{tikzpicture}[x=2.0cm, y=1.5cm, >=stealth, every node/.style={font=\small}]
    % --- 1. 領域の塗りつぶし (0 <= b <= 2a かつ b <= -2a + 2*sqrt(2a), a > 0) ---
    % b = 2a と b = -2a + 2*sqrt(2a) の交点は (a, b) = (0.5, 1)
    \fill[pattern=north east lines, pattern color=blue!65!black]
    (0, 0) -- (0.5, 1) -- plot[domain=0.5:2, samples=60] (\x, {2*sqrt(2*\x) - 2*\x}) -- cycle;

    % --- 2. 座標軸への垂線（点線） ---
    \draw[densely dotted, thick, gray!70!black] (0.5, 0) -- (0.5, 1) -- (0, 1);

    % --- 3. 座標軸 ---
    \draw[->, thick, gray!80!black] (-0.35, 0) -- (2.8, 0) node[right] {$a$};
    \draw[->, thick, gray!80!black] (0, -0.45) -- (0, 2.0) node[above] {$b$};
    \node[below left=2pt] at (0, 0) {$O$};

    % --- 4. 境界線の描画 ---
    % b = 2a（境界部分は実線、延長は点線）
    \draw[line width=1.4pt, blue!85!black] (0, 0) -- (0.5, 1);
    \draw[densely dotted, thick, blue!85!black] (0.5, 1) -- (0.85, 1.7) node[above] {$b = 2a$};
    % b = 0 上の境界 (0 < a <= 2)
    \draw[line width=1.4pt, blue!85!black] (0, 0) -- (2, 0);
    % 2a + b = 2*sqrt(2a)（領域外の 0<=a<=0.5 は点線、境界 0.5<=a<=2 は実線）
    \draw[densely dotted, thick, blue!85!black, domain=0:0.5, samples=30]
    plot (\x, {2*sqrt(2*\x) - 2*\x});
    \draw[line width=1.4pt, blue!85!black, domain=0.5:2, samples=60]
    plot (\x, {2*sqrt(2*\x) - 2*\x});
    \node[above right=1pt, blue!85!black, fill=white, inner sep=1pt, font=\footnotesize] at (0.72, {2*sqrt(1.44)-1.44}) {$2a+b = 2\sqrt{2a}$};

    % --- 5. 直線 2a + b = k ---
    \draw[densely dotted, thick, orange!90!black, domain=0.15:1.18]
    plot (\x, {2 - 2*\x}) node[below right, font=\footnotesize] {$2a+b=k$};
    \draw[line width=1.2pt, red!85!black, domain=1.65:2.22]
    plot (\x, {4 - 2*\x});
    \node[right=2pt, red!85!black, font=\footnotesize] at (2.15, -0.3) {$2a+b=4$};

    % --- 6. 軸目盛り・特徴点 ---
    \node[below=2pt, font=\scriptsize] at (0.5, 0) {$\frac{1}{2}$};
    \node[left=2pt, font=\scriptsize] at (0, 1) {$1$};
    \node[circle, fill, inner sep=1.2pt] at (0.5, 1) {};
    \node[circle, fill=red!85!black, inner sep=1.6pt] at (2, 0) {};
    \node[red!85!black, below=2pt] at (2, 0) {$2$};
    % a = 0 は除く
    \node[circle, draw=blue!85!black, thick, fill=white, inner sep=1.5pt] at (0, 0) {};
  \end{tikzpicture}
  \caption{$1^\circ$ における $(a,b)$ のとりうる領域（境界は $a=0$ のみ除く）}
\end{figure}

$f'(1)=2a+b$のとりうる範囲は直線$2a+b=k$と斜線部の領域の交点だから
\begin{align}
  0\le2a+b\le4
\end{align}
である．右の等号成立は$(a,b)=(2,0)$の時である．

\subparagraph{$2^{\circ}$の時（$0\le-\dfrac{b}{2a}\le 1$）}

$a>0$から$0 \le -b \le 2a$すなわち$-2a \le b \le 0$である．$|f(x)| \le 1$より$-1 \le \min f, \max f \le 1$であるから
\begin{align}
  \begin{dcases}
    -1 \le -\dfrac{b^2}{4a}+c \le 1 \\
    -1 \le a-b+c \le 1
  \end{dcases}
\end{align}
である．辺々引き算して
\begin{align}
  & -2 \le a-b+\dfrac{b^2}{4a} \le 2 \\
  & -8a \le 4a^2-4ab+b^2 \le 8a \\
  & -8a \le (2a-b)^2 \le 8a \\
  &\left|2a-b\right| \le 2\sqrt{2a} \\
  \therefore
  & 2a - 2\sqrt{2a} \le b \le 2a + 2\sqrt{2a}
\end{align}
である．これと$-2a \le b \le 0$を合わせて$(a,b)$平面に図示すると以下の斜線部である．

\begin{figure}[htb]
  \centering
  \begin{tikzpicture}[x=2.0cm, y=1.5cm, >=stealth, every node/.style={font=\small}]
    % --- 1. 領域の塗りつぶし (-2a <= b <= 0 かつ b >= 2a - 2*sqrt(2a), a > 0) ---
    % b = -2a と b = 2a - 2*sqrt(2a) の交点は (a, b) = (0.5, -1)
    \fill[pattern=north east lines, pattern color=blue!65!black]
    (0, 0) -- (0.5, -1) -- plot[domain=0.5:2, samples=60] (\x, {2*\x - 2*sqrt(2*\x)}) -- cycle;

    % --- 2. 座標軸への垂線（点線） ---
    \draw[densely dotted, thick, gray!70!black] (0.5, 0) -- (0.5, -1) -- (0, -1);

    % --- 3. 座標軸 ---
    \draw[->, thick, gray!80!black] (-0.35, 0) -- (2.8, 0) node[right] {$a$};
    \draw[->, thick, gray!80!black] (0, -1.9) -- (0, 0.65) node[above] {$b$};
    \node[above left=2pt] at (0, 0) {$O$};

    % --- 4. 境界線の描画 ---
    % b = -2a（境界部分は実線、延長は点線）
    \draw[line width=1.4pt, blue!85!black] (0, 0) -- (0.5, -1);
    \draw[densely dotted, thick, blue!85!black] (0.5, -1) -- (0.85, -1.7) node[below] {$b = -2a$};
    % b = 0 上の境界 (0 < a <= 2)
    \draw[line width=1.4pt, blue!85!black] (0, 0) -- (2, 0);
    % 2a - b = 2*sqrt(2a)（領域外の 0<=a<=0.5 は点線、境界 0.5<=a<=2 は実線）
    \draw[densely dotted, thick, blue!85!black, domain=0:0.5, samples=30]
    plot (\x, {2*\x - 2*sqrt(2*\x)});
    \draw[line width=1.4pt, blue!85!black, domain=0.5:2, samples=60]
    plot (\x, {2*\x - 2*sqrt(2*\x)});
    \node[below right=2pt, blue!85!black, fill=white, inner sep=1pt, font=\footnotesize] at (1.15, {2.3-2*sqrt(2.3)}) {$2a-b = 2\sqrt{2a}$};

    % --- 5. 直線 2a + b = k ---
    \draw[densely dotted, thick, orange!90!black, domain=0.45:1.22]
    plot (\x, {1.6 - 2*\x});
    \node[above right=1pt, orange!90!black, font=\footnotesize] at (0.48, 0.64) {$2a+b=k$};
    \draw[line width=1.2pt, red!85!black, domain=1.65:2.25]
    plot (\x, {4 - 2*\x});
    \node[right=2pt, red!85!black, font=\footnotesize] at (2.15, -0.3) {$2a+b=4$};

    % --- 6. 軸目盛り・特徴点 ---
    \node[above=2pt, font=\scriptsize, fill=white, inner sep=0.8pt] at (0.5, 0) {$\frac{1}{2}$};
    \node[left=2pt, font=\scriptsize] at (0, -1) {$-1$};
    \node[circle, fill, inner sep=1.2pt] at (0.5, -1) {};
    \node[circle, fill=red!85!black, inner sep=1.6pt] at (2, 0) {};
    \node[red!85!black, above=2pt] at (2, 0) {$2$};
    % a = 0 は除く
    \node[circle, draw=blue!85!black, thick, fill=white, inner sep=1.5pt] at (0, 0) {};
  \end{tikzpicture}
  \caption{$2^\circ$ における $(a,b)$ のとりうる領域（境界は $a=0$ のみ除く）}
\end{figure}

$f'(1)=2a+b$のとりうる範囲は直線$2a+b=k$と斜線部の領域の交点だから
\begin{align}
  0\le2a+b\le4
\end{align}
である．右の等号成立は$(a,b)=(2,0)$の時である．

\subparagraph{$3^{\circ}$の時}

この時は軸の条件から，
\begin{align}
  -\dfrac{b}{2a}\le-1 \ \text{or}\ 1\le-\dfrac{b}{2a}
\end{align}
である．$a>0$だから
\begin{align}
  2a\le b \ \text{or}\ b\le-2a \label{eq:9}
\end{align}
である．又$|f(\pm 1)| \le 1$より
\begin{align}
  & -1 \le a+ b + c \le 1 \\
  & -1 \le a- b + c \le 1
\end{align}
である．辺々足し引きして，
\begin{align}
  -1 \le b \le 1 \label{eq:10}
\end{align}
である．\cref{eq:9,eq:10}を図示すると下図斜線部である（境界は$a=0$のみ含まず）．

\begin{figure}[htb]
  \centering
  \begin{tikzpicture}[x=2.2cm, y=1.6cm, >=stealth, every node/.style={font=\small}]
    % --- 1. 領域の描画 (ハッチング) ---
    % (i) b >= 2a かつ b <= 1 (a > 0)
    \fill[pattern=north east lines, pattern color=blue!65!black]
    (0, 0) -- (0.5, 1.0) -- (0, 1.0) -- cycle;
    % (ii) b <= -2a かつ b >= -1 (a > 0)
    \fill[pattern=north east lines, pattern color=blue!65!black]
    (0, 0) -- (0.5, -1.0) -- (0, -1.0) -- cycle;

    % --- 2. 帯状領域 -1 <= b <= 1 の補助線 ---
    \draw[densely dotted, thick, gray!70!black] (0, 1.0) -- (1.45, 1.0) node[right, black] {$b=1$};
    \draw[densely dotted, thick, gray!70!black] (0, -1.0) -- (1.45, -1.0) node[right, black] {$b=-1$};
    \draw[densely dotted, thick, gray!70!black] (0.5, -1.0) -- (0.5, 1.0);

    % --- 3. 座標軸 ---
    \draw[->, thick, gray!80!black] (-0.35, 0) -- (2.1, 0) node[right] {$a$};
    \draw[->, thick, gray!80!black] (0, -1.75) -- (0, 1.85) node[above] {$b$};
    \node[below left=2pt] at (0, 0) {$O$};
    \node[left=2pt, font=\scriptsize] at (0, 1.0) {$1$};
    \node[left=2pt, font=\scriptsize] at (0, -1.0) {$-1$};
    \node[below=2pt, font=\scriptsize, fill=white, inner sep=1pt] at (0.5, 0) {$\frac{1}{2}$};

    % --- 4. 境界線 b = ±2a の描画 ---
    \draw[line width=1.4pt, blue!85!black] (0, 0) -- (0.5, 1.0);
    \draw[densely dotted, thick, blue!85!black] (0.5, 1.0) -- (0.78, 1.56) node[above left] {$b = 2a$};
    \draw[line width=1.4pt, blue!85!black] (0, 0) -- (0.5, -1.0);
    \draw[densely dotted, thick, blue!85!black] (0.5, -1.0) -- (0.78, -1.56) node[below left] {$b = -2a$};
    % 上下水平境界線 (0 < a <= 0.5)
    \draw[line width=1.4pt, blue!85!black] (0, 1.0) -- (0.5, 1.0);
    \draw[line width=1.4pt, blue!85!black] (0, -1.0) -- (0.5, -1.0);

    % --- 5. 直線 2a + b = k の描画（(0.5, 1) で最大値 k = 2） ---
    \draw[line width=1.2pt, red!85!black, domain=0.18:1.3]
    plot (\x, {2 - 2*\x}) node[below right, red!85!black, font=\footnotesize] {$2a+b=k$};

    % --- 6. 点のプロット（a = 0 上は除く白抜き丸） ---
    \node[circle, fill=red!85!black, inner sep=1.6pt] at (0.5, 1.0) {};
    \node[red!85!black, above right=2pt, fill=white, inner sep=1pt, font=\footnotesize] at (0.5, 1.0) {$\left(\frac{1}{2}, 1\right)$};
    \node[circle, draw=blue!85!black, thick, fill=white, inner sep=1.5pt] at (0, 0) {};
    \node[circle, draw=blue!85!black, thick, fill=white, inner sep=1.5pt] at (0, 1.0) {};
    \node[circle, draw=blue!85!black, thick, fill=white, inner sep=1.5pt] at (0, -1.0) {};
  \end{tikzpicture}
  \caption{$3^\circ$ における $(a,b)$ のとりうる領域（境界は $a=0$ のみ除く）}
\end{figure}

従って$f'(1)=2a+b$のとりうる範囲は，直線$2a+b=k$と斜線部の領域の交点を考えて
\begin{align}
  -1<2a+b\le2
\end{align}
だから$|2a+b|<4$である．
\casesend

以上から，$|2a+b|\le4$は成立する．

\paragraph{(2)}
(1)で等号が成立する場合を考えればよく，$1^{\circ}$および$2^{\circ}$と$a<0$の対称性（$f \mapsto -f$）も合わせて$(a,b)=(\pm 2,0)$の時である．（以下略）

\newpage

\end{document}
